\section{Application}

To demonstrate \mkado at genome scales,
we assembled a dataset of 12,437 human protein-coding genes
with both polymorphism and divergence data.
Codon-aligned orthologous sequences
for human (\textit{Homo sapiens}),
chimpanzee (\textit{Pan troglodytes}),
and orangutan (\textit{Pongo abelii})
were obtained from \texttt{OrthoMaM} v12 \citep{orthomam2024},
which provides one canonical coding-sequence ortholog per gene
across mammalian taxa.
Human intraspecific polymorphism was drawn from 178 Yoruba (YRI)
individuals in the 1000 Genomes Project NYGC high-coverage dataset
\citep{1000genomes2015},
yielding 356 phased haplotypes per gene.

\revpoint{1}{7}For each gene, polymorphism was projected onto the
\texttt{OrthoMaM} alignment in five steps:
(i)~the canonical Ensembl CDS for the gene was reconstructed from
the GRCh38 reference and the GTF-derived exon set, providing the
mapping from CDS position to genomic coordinate;
(ii)~the ungapped human row of the \texttt{OrthoMaM} alignment was
matched against the Ensembl CDS to obtain the alignment-column to
CDS-position correspondence (genes for which this match failed,
typically because of curation differences between
\texttt{OrthoMaM} and Ensembl, were dropped);
(iii)~for every CDS position that survived this match, the
overlapping YRI variants were extracted from the 1000 Genomes
NYGC VCF, with reference-allele identity confirmed and
minus-strand alleles complemented to match the CDS orientation;
(iv)~per-individual phased haplotypes were assembled from the
extracted variants, yielding 356 human haplotype sequences over
the Ensembl CDS positions retained by \texttt{OrthoMaM}; and
(v)~these haplotypes were re-gapped to the \texttt{OrthoMaM}
alignment, dropping any human-CDS columns absent from the
\texttt{OrthoMaM} alignment and inserting the alignment's
existing gap structure (so insertions, deletions, and
non-aligned regions in the multi-species alignment are
preserved). Chimpanzee and orangutan sequences were taken
directly from the \texttt{OrthoMaM} alignment. Multi-transcript
genes are reduced to the single canonical CDS chosen by the
\texttt{OrthoMaM} curation; alternative transcripts are not
analysed separately. The resulting per-gene multi-species FASTAs
are codon-aligned and ready for direct input to \mkado.

Using orangutan as the outgroup,
genome-wide pooled counts were
$D_n$\,=\,94,578, $D_s$\,=\,188,325, $P_n$\,=\,44,850,
and $P_s$\,=\,47,838.
The per-gene standard MK test identified numerous genes
with significant departures from neutrality,
including genes significant after Bonferroni correction
(volcano plot, left panel of \cref{fig:volcano})\revpoint{2}{3}.
As expected,
the majority of testable genes showed NI~$>$~1,
reflecting the well-characterized excess
of slightly deleterious nonsynonymous polymorphism in humans
\citep{li2008}.
Using chimpanzee as the outgroup yielded fewer fixed differences
($D_n$\,=\,37,120, $D_s$\,=\,65,228), roughly 2.5-fold less than
orangutan (left panel of \cref{fig:chimp_mk}),
and no individual genes survived Bonferroni correction,
consistent with the limited power of single-gene tests
at shorter divergence times.

The Tarone-Greenland weighted estimate \citep{stoletzki2011}
across all genes was strongly negative for both outgroups:
$\hat{\alpha}_\text{TG} = -0.853$
(95\% CI: $-0.888$ to $-0.812$)
for orangutan
and $\hat{\alpha}_\text{TG} = -0.651$
(95\% CI: $-0.682$ to $-0.607$)
for chimpanzee,
with corresponding $\text{NI}_\text{TG}$ values of 1.853 and 1.651.
These negative values confirm the pervasive excess
of weakly deleterious nonsynonymous polymorphism in humans.
Excluding singletons shifted both estimates upward
($\hat{\alpha}_\text{TG} = -0.794$ for orangutan;
$-0.606$ for chimpanzee),
as expected given the enrichment of slightly deleterious
variants at low frequency.
Applying a 15\% derived allele frequency cutoff
\citep{fay2001}
brought both estimates close to zero
($\hat{\alpha}_\text{TG} = -0.091$ for orangutan;
$-0.041$ for chimpanzee),
removing the bulk of low-frequency deleterious variants
but still failing to recover a signal of positive selection.
Similarly, the aggregated imputed MK test \citep{murga-moreno2022},
using a 15\% derived allele frequency cutoff,
estimated $\hat{\alpha} = -0.097$ for orangutan
and $\hat{\alpha} = -0.068$ for chimpanzee,
substantially reducing the downward bias
but likewise not recovering a positive signal of adaptation.

In contrast, the aggregated asymptotic MK test,
which models the frequency dependence of $\alpha$
and extrapolates to fixation \citep{messer2013},
successfully recovered a positive signal of adaptive evolution,
estimating $\hat{\alpha} = 0.163$
(95\% CI: 0.075 to 0.253; right panel of \cref{fig:volcano})
using orangutan as the outgroup,
indicating that approximately 16\% of amino acid substitutions
on the human lineage were driven by positive selection.
With chimpanzee as the outgroup the estimate widened to
$\hat{\alpha} = 0.075$ (95\% CI: $-1.22$ to $0.11$;
right panel of \cref{fig:chimp_mk}),
illustrating the greater uncertainty
from the shallower divergence, but also indicating 
a potential biological signal of less frequent adaptive evolution
after the human-Pongo split. 
These results are consistent with previous analyses
of human coding sequences \citep{messer2013}.

The entire analysis---standard MK tests,
Tarone-Greenland estimation, asymptotic estimation,
imputed estimation,
and volcano plot generation for all 12,437 genes---completed
in under one minute on a multi-core workstation
using \mkado's parallel batch mode.
